\subsection{Change of Base Field}

Let \(\mathrm{L}\) be an extension of the field \(\mathrm{K}\) . Let \(\mathrm{A}\) and \(\mathrm{B}\) be K-algebras of finite degree; then the L-algebras \(\mathrm{A}_{(\mathrm{L})}\) and \(\mathrm{B}_{(\mathrm{L})}\) are of finite degree. The L-algebras \(\mathrm{A}_{(\mathrm{L})} \otimes_{\mathrm{L}} \mathrm{B}_{(\mathrm{L})}\) and \((\mathrm{A} \otimes_{\mathrm{K}} \mathrm{B})_{(\mathrm{L})}\) are isomorphic (III, §4, No.~1, p.~462, Proposition 3). The L-algebra \(\mathrm{K}_{(\mathrm{L})}\) is isomorphic to \(\mathrm{L}\) . Consequently, there exists a unique monoid homomorphism \(\rho_{\mathrm{L/K}}\) from \(\mathcal{C}_{\mathrm{K}}\) to \(\mathcal{C}_{\mathrm{L}}\) such that

\[
\rho_ {\mathrm{L/K}} (\operatorname{cl} (\mathrm{A})) = \operatorname{cl} (\mathrm{A} _ {(\mathrm{L})})\tag{3}
\]

for every K-algebra A of finite degree.

If the K-algebras A and B are Morita equivalent, then so are the L-algebras \(A_{(L)}\) and \(B_{(L)}\) (VIII, p.~111, Proposition 13, e)). If the K-algebra A is central simple and of finite degree, then so is the L-algebra \(A_{(L)}\) (VIII, p.~251, Remark 2). So we can deduce from \(\rho_{L/K}\) a group homomorphism \(r_{L/K}\) from \(\mathrm{Br}(K)\) to \(\mathrm{Br}(L)\) such that

\[
r _ {\mathrm{L/K}} ([ \mathrm{A} ]) = [ \mathrm{A} _ {(\mathrm{L})} ]\tag{4}
\]

for every central simple K-algebra A of finite degree.

Let \(M\) be an extension of the field \(L\) . Because extension of scalars is transitive (III, §1, No.~5, p.~434), we have the relation

\[
r _ {\mathrm{M/K}} = r _ {\mathrm{M/L}} \circ r _ {\mathrm{L/K}}.\tag{5}
\]

Let A be a central simple K-algebra of finite degree, and let L be an extension of K. We say that L is a splitting field for A (or splits A) if the L-algebra \(\mathrm{A}_{(\mathrm{L})}\) is isomorphic to a matrix algebra \(\mathbf{M}_{n}(\mathrm{L})\) for an integer \(n \geqslant 1\) . In the aforementioned notation, this corresponds to saying that the class of A in \(\mathrm{Br}(\mathrm{K})\) belongs to the kernel of the homomorphism \(r_{\mathrm{L/K}} : \mathrm{Br}(\mathrm{K}) \to \mathrm{Br}(\mathrm{L})\) .

If B is similar to A, then L is a splitting field for A if and only if it is a splitting field for B.

If L is a splitting field for A, then every extension of L is a splitting field for A. By Theorem 1 of VIII, p.~252, there exists a Galois extension of K of finite degree that is a splitting field for A, and every separable closure of K is a splitting field for A.

PROPOSITION 5. --- Let A be a central simple K-algebra of finite degree, and let L be an extension of K of finite degree. The following properties are equivalent:

\begin{enumerate}
\def\labelenumi{(\roman{enumi})}
\item
  The extension L is a splitting field for A.
\item
  There exists a central simple K-algebra of finite degree similar to A containing a maximal commutative subalgebra isomorphic to L.
\end{enumerate}

Let us prove that (ii) implies (i); it suffices to consider the case when L is a maximal commutative subalgebra of A. Let \(\psi: A \otimes_{K} A^{\circ} \to \operatorname{End}_{K}(A)\) be the canonical isomorphism that sends \(a \otimes a'\) to the K-linear mapping \(x \mapsto axa'\) (VIII, p.~252, Theorem 1). We view A as a right vector space over L; then \(\psi\) sends \(A \otimes_{K} L\) into the subspace \(\operatorname{End}_{L}(A)\) of \(\operatorname{End}_{K}(A)\) and induces, by restriction, an injective homomorphism of L-algebras \(\psi': A \otimes_{K} L \to \operatorname{End}_{L}(A)\) . Set n = {[}L : K{]}. By VIII, p.~262, Proposition 3, we have {[}A : L{]} = n, and therefore \([A \otimes_{K} L : L] = [A : K] = n^{2}\) and \([\operatorname{End}_{L}(A) : L] = n^{2}\) . Consequently, \(\psi'\) is an isomorphism, which proves (i).

The converse follows from Lemma 3 below.

Lemma 3. --- Let A be a central simple K-algebra of finite degree, and let L be an extension of K of finite degree that is a splitting field for A. Let V be a simple \(\mathrm{A}_{(\mathrm{L})}\) -module, so that the natural morphism \(\varphi: \mathrm{A}_{(\mathrm{L})} \to \mathrm{End}_{\mathrm{L}}(\mathrm{V})\) is an isomorphism. Let C be the ring \(\mathrm{End}_{\mathrm{A}}(\mathrm{V})\) . Then C is similar to \(A^{o}\) , and the image of \(L \otimes 1 \subset A_{(L)}\) is a maximal commutative subalgebra of C.

We identify A with a subring of \(\mathrm{A}_{(\mathrm{L})}\) . We view V as a K-vector space. The ring C is the commutant of \(\varphi(\mathrm{A})\) in \(\operatorname{End}_{\mathrm{K}}(\mathrm{V})\) . It is a central simple K-algebra of finite degree, and the homomorphism \(a \otimes c \mapsto ac\) from \(A \otimes_{K} C\) to \(\operatorname{End}_{\mathrm{K}}(\mathrm{V})\) is an isomorphism (VIII, p.~260, Theorem 6, a)). Consequently, the K-algebras A and \(C^{o}\) are similar (VIII, p.~279, Proposition 3).

Let \(L_{V}\) be the ring of homotheties of the L-vector space V; it is the commutant of \(\mathrm{End}_{\mathrm{L}}(\mathrm{V})\) in \(\mathrm{End}_{\mathrm{K}}(\mathrm{V})\) (VIII, p.~82, Corollary 1). Now, the K-algebra \(\mathrm{End}_{\mathrm{L}}(\mathrm{V})\) is generated by \(\varphi(\mathrm{A})\) and \(L_{V}\) ; therefore, in \(\mathrm{End}_{\mathrm{K}}(\mathrm{V})\) , we have

\[
\mathrm{L} _ {\mathrm{V}} = \mathrm{End} _ {\mathrm{L}} (\mathrm{V}) ^ {\prime} = \varphi (\mathrm{A}) ^ {\prime} \cap \mathrm{L} _ {\mathrm{V}} ^ {\prime} = \mathrm{C} \cap \mathrm{L} _ {\mathrm{V}} ^ {\prime},
\]

where for any subset B of \(\operatorname{End}_{\mathrm{K}}(\mathrm{V})\) , the commutant of B in \(\operatorname{End}_{\mathrm{K}}(\mathrm{V})\) is denoted by \(B'\) . So \(L_{V}\) is a maximal commutative subalgebra of C (VIII, p.~261, Lemma 3), and therefore also of \(C^{o}\) . The mapping \(\lambda \mapsto \lambda_{V}\) is a K-algebra isomorphism from L to \(L_{V}\) . This proves the lemma.

COROLLARY 1. --- Let A be a central simple K-algebra of finite degree, and let L be an extension of K of finite degree. Suppose \([A:K]=[L:K]^{2}\) . Then L is a splitting field for A if and only if A contains a subalgebra isomorphic to L.

Suppose that there exists a morphism \(\varphi\) from L to A. Let M be a maximal semisimple commutative subalgebra containing \(\varphi(\mathrm{L})\) . By Proposition 3 of

VIII, p.~262, we have \([A:K] = [M:K]^2\) , hence \([M:K] = [L:K]\) and \(M = \varphi(L)\) . By Proposition 5, the field \(L\) splits \(A\) . Conversely, suppose that the field \(L\) splits \(A\) ; then it is isomorphic to a maximal commutative subalgebra of a central simple K-algebra \(B\) similar to \(A\) (Proposition 5). We have \([B:K] = [L:K]^2\) (VIII, p.~262, Proposition 3), hence \([B:K] = [A:K]\) . Consequently, \(B\) is isomorphic to \(A\) (VIII, p.~280, Corollary). Corollary 1 follows.

COROLLARY 2. --- Let D be a field of finite degree over K with center K, and let L be an extension of K of finite degree that is a splitting field for D. The reduced degree of D divides {[}L : K{]}.

Denote the reduced degree of D (VIII, p.~253) by r; by definition, we have \([D:K]=r^{2}\) . By Proposition 5, there exists a central simple K-algebra B similar to D of which L is a maximal commutative subalgebra. Since B is isomorphic to a matrix algebra \(\mathbf{M}_{n}(\mathrm{D})\) (VIII, p.~278, Lemma 1), we have \([B:K]=n^{2}r^{2}\) and consequently \([L:K]=nr\) (VIII, p.~262, Proposition 3).

